\Needspace{18\baselineskip}
\section{What remains for the paper's concrete theorem}
\label{app:remaining-work}
The unbounded contextual sufficiency theorem in the concrete encoding is
proved: Section~\ref{sec:unbounded-theorem} discharges the full-UAM depth-transfer
obligation and proves both implications. No unfinished Lean proof obligation
remains in that construction. The following paper-correspondence audits remain;
they qualify the earlier claim that the paper's theorem was fully settled.

\begin{enumerate}
\item \textbf{Context closure and observations.} Compare the paper's contextual
refinement with \texttt{PaperRefines} and the concrete observation construction.
Explain the absence of plugging closure from the final Lean premises, and
establish the paper's plugging-closed instance with matching observation
quantification, or correct the encoding if necessary.
See Section~\ref{sec:context-correspondence-audit}.
\item \textbf{Substitution stability and refocusing.} Audit the correspondence
between Figure~9(f), Lemma~7.6, and the formal substitution/rootedness argument.
Determine why the formal theorem needs no explicit stability premise and
whether the encoded rule and rootedness justify that difference.
See Section~\ref{sec:stability-correspondence-audit}.
\end{enumerate}

Output closure, input safety, robustness, deterministic and refocusing
behaviour, and a reasonable cost triple remain semantic hypotheses on an
application. The copying UAM, language translations, and application-specific
equivalences were outside the selected general-theorem milestone; Proposition~12.2
is now an explicit extension target.

\clearpage
\subsection{Remaining work for Proposition 12.2}
The machine, all eleven template families and their contextual consequences
are proved. Full semantic weakening and the actual left program's allocation
and beta-opening prefix are complete. Sections~\ref{sec:proposition-result-normalization}
and~\ref{sec:proposition-final-segment} now complete right-hand normalization
and concrete parametricity to the unchanged right-hand translation, in arbitrary
finite environments and at every adequate allowance. The bridge from the
actual opened left graph to the normalized private closure remains unproved.

Previously struck-through tasks have been removed. Only tasks completed in
this milestone are struck through below; every unstruck task remains open.
Earlier achievements remain documented in the main commentary.
% Progress policy: delete the previous milestone's completedtask entries;
% strike only achievements newly completed in the current milestone.
\begin{enumerate}
\item \textbf{Stability, renaming and substitution.}
\begin{itemize}
\item Prove any further stability/copyability invariants required by the
remaining specimen and translation arguments.
\item Prove non-bijective renaming and graph-level substitution laws.
\end{itemize}
\item \textbf{Specimens and observations.}
Prove correspondence between the coordinate specimens and the paper's
explicitly focused graphs, unique-hole inventories and replacement tuples.
Preserve the proved execution and observation laws through any necessary
semantic corrections, and retain the zero-input qualification of Lemma~E.1.
\item \textbf{The final example and its graph normalization.}
The two ends of the concrete equivalence chain are proved; its source-normalization
bridge is still open.
\begin{itemize}
\item {\raggedright\completedtask{Normalize the unchanged constant-result translation to the schematic constant closure, preserving exact successful counts at every adequate depth and termination observations in arbitrary finite environments.}\par}
\item {\raggedright\completedtask{Instantiate concrete parametricity, retaining any certified private sharing tree, and compose it with normalization to the unchanged right-hand translation in arbitrary finite and named environments.}\par}
\item Normalize the actual opened translated left body and its retained
sharing/box structure into the Figure~22 parametricity source, preserving
its private instance and store.
\item Carry the original-left allocation/beta prefix and source-normalization
certificates to the required common allowances, compose the complete chain,
and obtain the public Proposition~12.2 theorem for arbitrary distinct-list
environments.
\end{itemize}
\end{enumerate}
No claim of completion is made until the final general-environment theorem
passes the build and axiom audit. The current estimate remains 70--80\% of
mathematical effort, with low confidence; the number of completed declarations
or checklist items is not a measure of the remaining difficulty.
